% !TEX program = lualatex
\documentclass[11pt,a4paper,ja=standard,autodetect-engine]{bxjsarticle}
\usepackage{amsmath,amssymb,amsthm}
\usepackage{bm}

\newtheorem{theorem}{定理}
\newtheorem{definition}{定義}
\newtheorem{proposition}{命題}
\newtheorem{example}{例}
\newtheorem{remark}{注意}

\title{継続残余としての証明障害：\\ 実行可能性と静的証明のあいだの初等モデル}
\author{千葉将臣}
\date{2026\mbox{-}06\mbox{-}29}

\begin{document}

\maketitle

\begin{abstract}
本論文は、動的な計算過程を古典的な静的証明へ写すときに現れる差分を、継続残余 $\rho$ として記述する初等的モデルを提案する。出発点は、直観主義論理では一般に二重否定除去 $\neg\neg A\to A$ が許されず、Curry\mbox{--}Howard 対応のもとで古典的制御は継続構造として解釈される、という標準的な観察である。本稿ではこの継続構造のうち、構成的な証拠として直接観測されない評価文脈を残余と呼ぶ。さらに、プログラムの \texttt{if} 分岐、メビウス帯の向き反転、コラッツ写像の軌道検証を例として、実行時には追跡できるが、実行前に静的な $\rho=0$ 証明として圧縮する際に残余指標が現れる、というモデルを与える。本稿の主張はコラッツ予想の ZFC 独立性ではなく、動的検証から静的証明への翻訳に生じる証明圧縮上の障害を可視化することである。
\end{abstract}

\section{導入}

数学的証明は通常、完成した静的対象として提示される。一方で、計算機上の検証や構成的な推論では、命題はしばしば実行可能な過程として現れる。たとえば、ある入力 $n$ に対してプログラムを実行し、有限時間後に望ましい出力が得られることを確認することはできる。しかし、その確認手順全体を実行前から一つの静的定理として固定することは、一般に別の問題である。

本稿の目的は、この差を「真偽の差」ではなく「翻訳の差」として扱うことである。すなわち、動的な実行プロトコル $P$ を古典的な証明対象へ写すとき、観測可能な部分 $P_{\mathrm{obs}}$ と、評価文脈として残る継続成分 $\mathcal{R}(P)$ とに分ける。そして、後者そのものを古典的対象として完全保存するのではなく、その影を残余指標 $\rho(P)$ として記録する。

この立場では、$\rho=0$ は「実行前に古典的証明として完全圧縮できる」ことを意味し、$\rho>0$ は「実行時の継続成分が静的証明へ完全には消去されていない」ことを意味する。ただし、$\rho>0$ は直ちに命題の独立性や反証を意味しない。これはあくまで、動的検証から静的証明へ移る際の証明圧縮上の残余である。

\section{論理的背景：二重否定と継続}

直観主義論理では、命題 $A$ に対して一般には
\begin{equation}
    \neg\neg A \to A
\end{equation}
を導出できない。これは、$A$ の否定が矛盾を導くことと、$A$ の構成的証拠を直接与えることが区別されるためである。

Curry\mbox{--}Howard 対応のもとでは、証明はプログラムに、命題は型に対応する。この対応において、古典論理の二重否定除去や排中律は、継続、例外、制御演算子のような計算構造として解釈される。つまり、古典的には $A$ と見なされる証明の内部に、構成的には「$A$ を否定する文脈を受け取って全体を制御する」評価文脈が含まれる。

本稿では、この評価文脈を残余の原型とみなす。

\begin{definition}[継続残余]
命題 $A$ の古典的証明を構成的・計算論的に読むとき、直接の値または構成的証拠として観測される部分を $A_{\mathrm{obs}}$ とし、二重否定除去、排中律、例外処理、または \texttt{if} 分岐の未選択枝として評価文脈に残る部分を $\mathcal{R}(A)$ とする。この $\mathcal{R}(A)$ の古典的に観測可能な指標を継続残余 $\rho(A)$ と呼ぶ。
\end{definition}

\begin{remark}
ここでいう「存在」は、直観主義論理と古典論理の差分を継続構造として読むという意味での存在である。したがって、任意の具体的命題について $\rho(A)>0$ が自動的に従うわけではない。各問題に対して、どの翻訳を採用し、どの継続成分を観測不能成分とみなすかを明示する必要がある。
\end{remark}

\section{残余翻訳}

動的な検証手順を $P$ とする。$P$ は、実行時には入力、分岐、ループ、例外処理を通じて値を返す。一方、古典的証明では、これらの実行過程を停止した静的対象として記述する必要がある。

\begin{definition}[残余翻訳]
動的証明プロトコル $P$ に対し、古典的観測者が静的対象として固定できる成分を $P_{\mathrm{obs}}$、固定しきれない継続成分を $\mathcal{R}(P)$ とする。残余翻訳とは、写像
\begin{equation}
    \mathcal{T}_{\rho}(P)=(P_{\mathrm{obs}},\rho(P))
\end{equation}
である。ここで $\rho(P)$ は $\mathcal{R}(P)$ そのものではなく、その分岐差分、ホロノミー欠損、または二重否定除去コストとして古典的に観測可能な不変量へ写したものである。
\end{definition}

この定義により、古典的証明と動的検証の差を次のように表せる。
\begin{itemize}
    \item $\rho(P)=0$：$P$ は実行前に静的証明へ完全圧縮される。
    \item $\rho(P)>0$：$P$ の継続成分の影が静的観測面に残る。
\end{itemize}

重要なのは、$\rho(P)>0$ が $P$ の失敗を意味しないことである。むしろ、$P$ は実行時には正しく動くが、その正しさを静的対象として固定する際に追加の指標が必要になる、という状況を表す。

\section{分岐差分としての残余}

最も初等的なモデルは、プログラムの \texttt{if} 分岐である。ある状態 $s$ に対して、プログラムが
\begin{equation}
    P(s)=
    \begin{cases}
        P_{1}(s) & \text{if } b(s)=1,\\
        P_{0}(s) & \text{if } b(s)=0
    \end{cases}
\end{equation}
として実行されるとする。

実行時には、$b(s)$ の値に応じて一方の枝だけが選ばれる。しかし静的証明では、選ばれなかった枝も含め、分岐全体が同一の証明対象へ収束することを示す必要がある。このとき、分岐差分を
\begin{equation}
    \partial_{\rho}P(s)=P_{1}(s)-P_{0}(s)
\end{equation}
と形式的に書く。実際には $P_{1}(s)$ と $P_{0}(s)$ が数値でない場合もあるため、これは「二つの枝の観測差」を表す記号である。

\begin{definition}[分岐残余]
分岐プログラム $P$ に対し、二つの枝が古典的観測面で同一視できるとき $\partial_{\rho}P=0$ と書く。同一視できない差分が残るとき、適当な大きさ関数 $\|\cdot\|$ により
\begin{equation}
    \rho(P)=\|\partial_{\rho}P\|
\end{equation}
と定める。
\end{definition}

このモデルでは、$\rho=0$ は「どちらの枝を通っても静的証明として同じ対象に戻る」ことを意味する。$\rho>0$ は「実行時には一方の枝が処理されるが、静的観測では未選択枝の影が残る」ことを意味する。

\section{メビウス帯：残余インデックスの最小例}

メビウス帯は、残余を直観的に説明する最小例である。帯の上を一周すると、局所的には連続的に進んでいるだけである。しかし始点へ戻ったとき、向きは反転している。二周すると元に戻る。

この構造を、状態集合 $X$ と符号 $\varepsilon\in\{+1,-1\}$ を持つ系として表す。メビウス輸送を
\begin{equation}
    M(x,\varepsilon)=(x',-\varepsilon)
\end{equation}
と書けば、二回の輸送では
\begin{equation}
    M^{2}(x,\varepsilon)=(x'',\varepsilon)
\end{equation}
となり、符号は回復する。

\begin{definition}[メビウス残余]
一周後の符号反転を古典的な単一平面上で消去しようとするとき、観測面に残る向きの差分を
\begin{equation}
    \rho_{\mathrm{Mob}}=1
\end{equation}
と記録する。二周後には反転が相殺されるため、合成された観測では
\begin{equation}
    \rho_{\mathrm{Mob}}^{(2)}=0
\end{equation}
と書く。
\end{definition}

この例が示すのは、動的な輸送そのものは矛盾なく実行できるが、一周分のねじれを単一の静的平面に完全消去しようとすると、向き反転のインデックスが残る、ということである。ここでの $\rho$ は矛盾や反証ではなく、継続された輸送の曲率的な痕跡である。

\section{コラッツ写像と実行前圧縮の残余}

コラッツ写像 $C:\mathbb{N}\to\mathbb{N}$ を
\begin{equation}
    C(n)=
    \begin{cases}
        n/2 & (n\equiv 0 \pmod 2),\\
        3n+1 & (n\equiv 1 \pmod 2)
    \end{cases}
\end{equation}
で定める。コラッツ予想は
\begin{equation}
    \mathrm{Col}:\quad \forall n\in\mathbb{N}\,\exists k\in\mathbb{N}\; C^{k}(n)=1
\end{equation}
という形を持つ。

各 $n$ を固定すれば、プログラムを実行して軌道を追跡できる。この意味で、有限範囲の追試はランタイム検証である。一方、静的証明では、全ての $n$ について停止時刻を統制する構造を実行前に与える必要がある。構成的には、これは停止時刻関数
\begin{equation}
    K:\mathbb{N}\to\mathbb{N},\qquad C^{K(n)}(n)=1
\end{equation}
を何らかの有限な証明原理によって支配することに近い。

\begin{definition}[コラッツ残余]
コラッツ検証プロトコル $P_{\mathrm{Col}}$ に対し、個々の入力で実行時に得られる停止確認を $P_{\mathrm{run}}(n)$ とする。これらを実行前に一つの静的証明対象へ圧縮する際に必要となる未回収の分岐・停止時刻・評価文脈の指標を
\begin{equation}
    \rho_{\mathrm{Col}}=\rho(P_{\mathrm{Col}})
\end{equation}
と呼ぶ。
\end{definition}

この定義のもとで、$\rho_{\mathrm{Col}}=0$ は、コラッツ軌道の全入力に対する停止確認が実行前に静的証明として完全圧縮されている状態を表す。$\rho_{\mathrm{Col}}>0$ は、個々の実行では停止を確認できても、その確認列全体を一つの静的証明対象へ移す段階で、分岐構造または停止時刻構造の残余が残ることを表す。

\begin{proposition}[実行検証と静的圧縮の非対称性]
有限個の入力に対するコラッツ軌道の実行検証は、コラッツ予想の静的証明とは異なる。前者は各入力に対するランタイム上の到達確認であり、後者は全入力を支配する実行前の圧縮原理を必要とする。
\end{proposition}

\begin{proof}
有限集合 $F\subset\mathbb{N}$ に対しては、各 $n\in F$ の軌道を計算し、ある $k_{n}$ について $C^{k_{n}}(n)=1$ を確認すればよい。しかし、コラッツ予想は全称命題であり、任意の $n$ に対して停止時刻の存在を示す必要がある。したがって、有限個の実行履歴の集積と、全入力を支配する静的証明とは論理的に区別される。
\end{proof}

\begin{remark}
本稿の主張は、コラッツ型の分岐実行を静的な $\rho=0$ 証明へ翻訳する際に、停止時刻関数と分岐構造に由来する残余指標を導入できる、というモデル的主張である（これはコラッツ予想が ZFC 独立であることを意味しない）。また、標準的なコラッツ予想が停止性問題と同値であるとも主張しない。
\end{remark}


\section{ZFC+U 語彙における $\rho_{\mathrm{Col}}=0$ 証明の未完結性}

前節までで、コラッツ写像そのものの真偽ではなく、コラッツ型分岐実行を静的な証明対象へ翻訳する際の残余指標 $\rho_{\mathrm{Col}}$ を導入した。本節では、この指標に関するモデル内命題を述べる。ここでいう未完結性は、コラッツ予想の ZFC 独立性を意味するものではなく、$\rho_{\mathrm{Col}}=0$ という完全圧縮条件を ZFC+U の語彙だけで閉じることの困難を表す。

\begin{proposition}[$\rho_{\mathrm{Col}}=0$ 証明の ZFC+U 語彙内未完結性]
本稿の残余翻訳モデルにおいて、$\rho_{\mathrm{Col}}=0$ を示すためには、停止時刻関数
\begin{equation}
    K:\mathbb{N}\to\mathbb{N},\qquad C^{K(n)}(n)=1
\end{equation}
を全入力に対して与えるだけでなく、その構成が全ての分岐継続を消去していることを示す必要がある。ZFC+U の通常の静的語彙は、この継続消去の完了を直接表す原理を含まない。したがって、$\rho_{\mathrm{Col}}=0$ の証明は、ZFC+U の語彙だけでは完結した対象として記述されず、追加の残余翻訳または継続消去原理を必要とする。
\end{proposition}

\begin{proof}
$\rho_{\mathrm{Col}}=0$ とは、個々の入力に対する実行時の停止確認が、未回収の分岐・停止時刻・評価文脈を残さず、一つの静的証明対象へ完全圧縮されることを意味する。コラッツ命題は
\begin{equation}
    \forall n\in\mathbb{N}\,\exists k\in\mathbb{N}\; C^{k}(n)=1
\end{equation}
という形を持つため、構成的に読むならば、各入力 $n$ に停止時刻 $K(n)$ を割り当てる全域的構成が必要になる。

しかし、本稿の残余翻訳では、$K$ の値を与えるだけでは十分ではない。$\rho_{\mathrm{Col}}=0$ は、各 $K(n)$ の存在に加えて、奇偶分岐によって発生する全ての未選択枝、評価文脈、停止時刻の探索過程が、静的証明対象の内部で完全に回収されていることを要求する。すなわち、必要なのは単なる停止時刻関数ではなく、停止時刻関数の構成が継続残余を残さないことの証明である。

ZFC+U の通常の語彙は、集合、関数、関係、量化によって静的対象を記述する。一方、ここで要求される「継続消去の完了」は、定義1の意味での $\mathcal{R}(P_{\mathrm{Col}})$ の完全消去である。すなわち、定義1の意味での $\mathcal{R}(P_{\mathrm{Col}})$ の完全消去は、集合・関数・関係・量化からなる ZFC+U の通常の語彙には属さない条件として扱われる。記号的には、
\begin{equation}
    \mathcal{R}(P_{\mathrm{Col}})\text{-Elim}
    \notin
    \mathrm{Lang}(\mathrm{ZFC+U})
\end{equation}
と表せる。ここで $\mathcal{R}(P_{\mathrm{Col}})\text{-Elim}$ は定義1の意味での継続成分の完全消去を表す略記であり、$\mathrm{Lang}(\mathrm{ZFC+U})$ は集合・関数・関係・量化による通常語彙を表す。したがって、この条件を ZFC+U 内の通常の対象として表すには、残余翻訳 $\mathcal{T}_{\rho}$ または同等の継続消去原理を追加で明示しなければならない。

よって、本稿のモデルにおいては、$\rho_{\mathrm{Col}}=0$ の証明は ZFC+U の語彙だけでは閉じず、少なくとも残余翻訳または継続消去を記述する追加語彙を必要とする。
\end{proof}

\begin{remark}
本命題は ZFC+U における証明可能性の独立性を主張しない。主張しているのは、$\rho_{\mathrm{Col}}=0$ という完全圧縮条件を記述するために必要な継続消去の語彙が ZFC+U に含まれないという、語彙の層の問題である。コラッツ予想そのものが ZFC+U で証明可能かどうかは、本稿の射程外である。ただし、$\rho_{\mathrm{Col}}=0$ がコラッツ予想の静的証明の完成と同値である以上、その記述に必要な語彙が ZFC+U に届かないという事実は、ZFC+U の語彙内でのコラッツ予想の完全な証明記述が原理的に困難であることを示唆する。
\end{remark}

\section{残余を古典的に見せるための翻訳コード}

残余そのものは継続構造であり、古典的な静的対象へそのまま戻すことは難しい。そこで、本稿では残余を次の三種類の古典的指標として翻訳する。

\begin{enumerate}
    \item \textbf{分岐差分}：\texttt{if} の二つの枝が静的観測面で同一対象へ戻るかを測る。
    \item \textbf{ホロノミー欠損}：ループを一周した後に、状態、型、符号、向きが恒等的に戻るかを測る。
    \item \textbf{二重否定除去コスト}：$\neg\neg A$ から $A$ を取り出すときに必要な継続制御の有無を測る。
\end{enumerate}

これらをまとめて、残余翻訳コードを
\begin{equation}
    \mathcal{T}_{\rho}(P)=
    \bigl(P_{\mathrm{obs}},
    \partial_{\rho}P,
    \mathrm{Hol}_{\rho}(P),
    \mathrm{Cost}_{\neg\neg}(P)\bigr)
\end{equation}
と書く。ここで $P_{\mathrm{obs}}$ は古典的に固定可能な証明成分、残りの三項は継続成分の影である。

この形式化の利点は、古典的観測者に対して「残余は見えない」と言うだけでなく、「残余そのものは見えないが、分岐差分、ホロノミー欠損、二重否定除去コストとしてその影は見える」と述べられる点にある。

\section{結論}

本稿では、直観主義論理と古典論理の差分を継続構造として読み、そのうち構成的証拠として直接観測されない部分を継続残余 $\rho$ と呼ぶ初等モデルを提示した。さらに、\texttt{if} 分岐、メビウス帯、コラッツ写像を通じて、動的実行では追跡できるが、実行前の静的証明へ圧縮する際に残余指標が現れる、という構図を説明した。

本稿の立場では、証明とは単に真なる命題を記号列として保存することではなく、動的検証プロトコルをどの程度まで静的対象へ圧縮できるかという翻訳問題でもある。$\rho=0$ は完全圧縮を、$\rho>0$ は継続成分の影が残ることを表す。今後の課題は、残余指標を具体的な形式体系内で符号化し、各問題に対して $\rho$ を比較可能な不変量として計算することである。

\end{document}
